\documentclass[10pt,a4paper]{amsart}
\usepackage[margin=19mm]{geometry}
\usepackage{amsmath,amssymb,mathtools}
\usepackage[scr=rsfs,cal=euler]{mathalfa}
\usepackage{xfrac}
\usepackage[unicode,hidelinks,bookmarks=false]{hyperref}
\newcommand{\ie}{i.e.\ }
\newcommand{\cf}{cf.\ }
\newcommand{\resp}{respectively\ }
\newcommand{\Subsection}{\subsection}
\providecommand{\nobreakdash}{\nobreak\hbox{-}}
\setlength{\emergencystretch}{2em}
\multlinegap0pt
\usepackage{amssymb}
\usepackage{xcolor}
\usepackage{mathtools}
\newtheorem{lemma}{Lemma}
\newcommand{\Z}{\mathbb{Z}}
\newcommand{\attop}[1]{{\let\textstyle\scriptstyle\let\scriptstyle\scriptscriptstyle\substack{#1}}}
\newcommand{\ints}{\cap}
\newcommand{\x}{\times}
\title{Bredon homological stability for configuration spaces of $G$-manifolds}
\author{Eva Belmont, J.D. Quigley,, Chase Vogeli}
\date{Source: arXiv:2311.02459v1 (CC BY 4.0)}
\begin{document}
\maketitle

\noindent\emph{Source and licence.} Bredon homological stability for configuration spaces of $G$-manifolds. Version arXiv:2311.02459v1 (CC BY 4.0), distributed by arXiv under the Creative Commons Attribution 4.0 International licence, http://creativecommons.org/licenses/by/4.0/, as recorded in the arXiv licence metadata for this exact version and shown on the abstract page https://arxiv.org/abs/2311.02459v1. Full licence notice: \texttt{licenses/arxiv-2311.02459-CC-BY-4.0.txt}.\par\smallskip

\noindent\textit{Editorial scope.} Counted unit: the article's lemma lem:fg-homology (source0.tex lines 1990--1993). The article's complete original proof of it is delivered with it (source0.tex lines 1994--2049, one \texttt{proof} environment of 56 lines). No mathematics is written, repaired or re-derived: the statement and the proof are the article's own lines, in the article's own notation, and no retained line is substituted or re-flowed.  Retained uncounted as the source-local context the proof needs: lines 1571--1574; lines 1934--1942; lines 1966--1971.  Nothing was omitted from any retained span, so the package carries no omission record.  No retained line contains a citation, so the wrapper carries no bibliography and no bibliography entry.  No retained line contains a figure environment, an image-inclusion command or drawing code, so no image is delivered and the package declares no asset dependency.  Editorial formatting only, in addition: an AMS \texttt{amsart} preamble that restates the article's own declarations and macros used by the retained lines, namely the environments \texttt{lemma} on separate counters, together with the standard packages those lines need.  The retained environments print in this document as Lemma 1 in order of appearance, whereas the article prints the same items with the numbering of its own full document.  The complete native source of the article as distributed in the arXiv e-print for this exact version is bundled unchanged as \texttt{sources/149543/source0.tex}.
\begin{lemma} \label{lem:homology-orbits-finite}
Suppose $X$ has a free $G$-action, and $H_d(X;\Z)$ is a finitely generated
abelian group for all $d$. Then $H_d(X/G;\Z)$ is finitely generated.
\end{lemma}

\begin{lemma} \label{lem:hypothesis-translation}
Suppose $G$ is Dedekind or a $p$-group, and let $W = W_GH$ denote the Weyl
group. If $M^H/W$ is connected
%and $M^H$ is an open (i.e., noncompact) submanifold of $M$,
then $M_{(H)}/G$ (if nonempty) is
%open and
connected. Moreover,
$$ M_{(H)}/G \cong \Big(M^H \backslash \bigcup_{e < L \leq W} (M^H)^L\Big)/W. $$
\end{lemma}

\begin{lemma}\label{lem:local-coeffs}
Suppose $M$ is $G$-manifold that is the interior of a compact $G$-manifold
with boundary. Let $H\leq G$. If $A$ is a $\pi_1(M^H)$-module that is finitely
generated as an abelian group, then the homology with local coefficients
$H_q(M^H;A)$ is a finitely generated abelian group for every $q$.
\end{lemma}

\begin{lemma} \label{lem:fg-homology}
Assume $G$ is abelian, and $M$ is the interior of a compact $G$-manifold.
Then $H_d(M_{(H)}/G; \Z)$ is finitely generated for all $d\in \Z,H\leq G$.
\end{lemma}

\begin{proof}
Let $W = W_GH$ be the Weyl group.
Lemma \ref{lem:hypothesis-translation} says that $M_{(H)}/G \cong X/W$ where
\begin{equation}\label{eq:bigcupX} X = M^H \backslash \bigcup_{e < L\leq W} (M^H)^L.
\end{equation}
By Lemma \ref{lem:homology-orbits-finite}, it suffices to show that
$H_*(X;\Z)$ is finitely generated in each
degree.

If $q:G\to W = G/H$ is the quotient, then $(M^H)^L = M^{q^{-1}(L)}$.
Abusing notation, we write $M^L = (M^H)^L$; also let
$U_L = M^H \backslash M^L$.
Since $M^L$ is a submanifold, there is a normal bundle $N(M^L)$ which is open,
and $N(M^L) \cup U_L \cong M^H$.
First we will show $H_d(U_L)$ is finite using a Mayer--Vietoris sequence.

The intersection term $N(M^L)\cap U_L$ is homotopy equivalent to the sphere bundle $S(N(M^L))$. There
is a Serre spectral sequence associated to the fibration $S^n\to S(N(M^L))\to
M^L$,
$$ E_2^{*,*} = H_*(M^L; H_*(S^n; \Z))\Rightarrow H_*(S(N(M^L));\Z) $$
where $\pi_1(M^L)$ may act nontrivially on $H_*(S^n;\Z)$. By Lemma
\ref{lem:local-coeffs}, the $E_2$ term is finitely generated in each degree, and
hence so is the target.

Now consider the Mayer--Vietoris sequence
$$ \dots\to H_{d+1}(M^H)\to H_d(S(N(M^L))) \to H_d(N(M^L)) \oplus H_d(U_L)\to
\dots. $$
The first term is finitely generated by Lemma \ref{lem:local-coeffs} and above we showed the second is
as well. Hence the third term is finitely generated; in particular so is $H_d(U_L)$.

Since $G$ is abelian, write $W = \prod_i \Z/p_i^{r_i}$.	In the union \eqref{eq:bigcupX}, it
suffices to restrict $L$ to products of elementary abelian subgroups
$\Z/p_i\cong \langle p_i^{r_i-1}\rangle\leq \Z/p_i^{r_i}$, since any nontrivial
subgroup $H\leq \Z/p_i^{r_i}$
contains this $\Z/p_i$. If $L\x L' \leq W$, we have $(M^H)^L \ints (M^H)^{L'} =
(M^H)^{L\x L'}$. More generally,
given $L = \prod_i L_i$ and $L' = \prod_i L'_i$, note
that $LL'\leq W$ is a product of all $L_i$ and $L'_i$ factors, with duplicates
removed. Thus we have
$$ U_L \cup U_{L'} = M^H \backslash ((M^H)^L \cap (M^H)^{L'}) = M^H \backslash
(M^H)^{LL'}. $$
We will show that 
$$X = \bigcap_\attop{L = L_1 \x \dots \x L_n\\L_i\text{ elem.ab.}} U_L $$
has finitely generated homology, by induction on the number of sets intersected.
The base case was shown above. Suppose any intersection of $n$ such sets has
finitely generated homology, and consider $U_{L_1}\ints \dots \ints U_{L_n}
\ints U_L$. There is a Mayer-Vietoris sequence
$$ \dots \to H_{d+1}((\cap_{i=1}^n U_{L_i})\cup U_L) \to H_d((\cap_{i=1}^n
U_{L_i})\cap U_L) \to H_d(\cap_{i=1}^n U_{L_i}) \oplus H_d(U_L) \to \dots. $$
The inductive hypothesis implies that the third term is finitely generated, so
in order to show the second term is finitely generated, it suffices to show the first
term is finitely generated. We have
$$ \Big(\bigcap_{i=1}^n U_{L_i}\Big) \cup U_L = \bigcap_{i=1}^n (U_{L_i} \cup
U_L) = \bigcap_{i=1}^n U_{L_i L} $$
and the inductive hypothesis shows this has finitely generated homology.
\end{proof}

\end{document}
